% !TeX root = ../main.tex


Large language models (LLMs) are increasingly used for code-related tasks, such as code generation, completion, and test generation. These tasks often incur substantial token costs because conventional programming languages are designed primarily for human reading and writing, and consequently contain syntactic, formatting, and structural information that may be redundant for LLM processing. Existing approaches reduce this overhead by heuristically defining concise syntax rules or mining frequently occurring patterns. However, frequency alone does not indicate whether a pattern is suitable for compression, and indiscriminate application of frequency-based rules can degrade downstream performance. We present \tool, a reversible code representation designed to reduce the token overhead of code language models. \tool organizes code compression into a unified \textit{Normalization \& Hybrid Idioms Shorthanding} framework. First, \textit{Normalization} leverages Python grammar and decoder state to normalize layout and structural information that can be deterministically recovered during decoding. Then, \textit{Hybrid Idioms Shorthanding} mines recurring code idioms at multiple granularities from the normalized representation and encodes them using compact symbols with typed captures. To avoid performance degradation from indiscriminate compression, \tool further introduces \textit{Conditional Attention}, which uses an offline model-derived signal to assess each eligible shorthand occurrence and selectively materialize it in compressed form, thereby incorporating model-level information into the compression process. Our evaluation shows that \tool reduces token consumption by up to 26.44\% among the evaluated modern code tokenizers. Downstream experiments demonstrate that code models can learn, generate, and use the \tool dialect, retaining performance close to the original Python surface in most evaluated backbone--task combinations while exposing a model-dependent gap on Mistral generation. Ablation and sensitivity studies further demonstrate the contributions of the individual components and the robustness of the selected parameters.
